\begin{wraptable}{r}{0.48\textwidth}
% \setlength{\tabcolsep}{1.5mm}
% \renewcommand{\arraystretch}{1.15}
\centering
\vspace{-17pt}
\caption{
\textbf{Estudio de ablación sobre \textsc{QaEgo4D}$_\texttt{test-mc}$}.
\colorbox[RGB]{235,235,235}{``Oracle Retrieval''} se refiere a un escenario en el que los segmentos de video anotados y relevantes para la pregunta se usan como entrada, con un uniform sampling de hasta 16 frames.
Esta configuración, por definición, tiene 100\% de recall y define el límite superior del rendimiento de VideoQA.
}
\label{tab:ablation_qaego4d}
\vspace{-4pt}
\resizebox{\linewidth}{!}{
\begin{tabular}{l cc}
    \toprule
    Método de retrieval    & VideoQA Acc. & Recall \\
    \midrule
    \footnotesize{\texttt{\textcolor{gray}{LLaVA-OV-0.5B}}} \\
    Uniform Sampling    & 42.6    & 6.1 \\
    External Retrieval  & 48.0    & 58.1 \\
    Internal Retrieval  & 50.0    & 63.4 \\
    \rowcolor[RGB]{235,235,235}
    Oracle Retrieval    & 52.0    & 100 \\
    \midrule
    \footnotesize{\texttt{\textcolor{gray}{LLaVA-OV-7B}}} \\
    Uniform Sampling    & 53.0    & 6.1 \\
    External Retrieval  & 54.2    & 58.1 \\
    Internal Retrieval  & 56.0    & 70.5\\
    \rowcolor[RGB]{235,235,235}
    Oracle Retrieval    & 64.4    & 100 \\
    \bottomrule
\end{tabular}
}
\vspace{-15pt}
\end{wraptable}
